% =====================================================================
%  CIKM submission --- anonymous review copy.
%  Build:  latexmk -pdf main.tex      (or: pdflatex -> bibtex -> pdflatex x2)
%  All numbers in tables/figures are generated from recorded benchmark
%  artefacts by make_assets.py; nothing is hand-transcribed.
% =====================================================================
\documentclass[sigconf, review, anonymous]{acmart}

% --- review-time boilerplate suppression -----------------------------
\settopmatter{printacmref=false}
\settopmatter{printfolios=true}
\renewcommand\footnotetextcopyrightpermission[1]{}
\acmConference[Anonymous submission]{}{}{}
\acmISBN{}
\acmDOI{}

% --- margin line numbers -------------------------------------------------
% The `review' class option above makes acmart print a per-line counter down
% BOTH margins and hard-codes its colour to red (acmart.cls \ACM@mk@linecount,
% \ACM@linecountL/R).  That is a review convenience supplied by the template,
% not a CIKM requirement, and a red-numbered page must never ship as the
% camera-ready.  Default off: define \reviewlinestrue before \begin{document}
% only if a venue explicitly requests a line-numbered copy.
\newif\ifreviewlines
\reviewlinesfalse
\makeatletter
\ifreviewlines\else
  \def\ACM@linecountL{}%
  \def\ACM@linecountR{}%
\fi
\makeatother

\usepackage{booktabs}
\usepackage{graphicx}
\usepackage{subcaption}
% acmart already defines \Bbbk; amssymb would clash. Release it first.
\let\Bbbk\relax
\usepackage{amsmath,amssymb}
\usepackage{algorithm}
\usepackage{algpseudocode}
\usepackage{xcolor}
\usepackage{enumitem}
\usepackage{balance}
\usepackage{pgfplots}
\pgfplotsset{compat=1.17}
\usetikzlibrary{arrows.meta,positioning,fit,shapes.geometric,calc}

% compact lists
\setlist[itemize]{leftmargin=1.1em,topsep=2pt,itemsep=1pt,parsep=0pt}
\setlist[enumerate]{leftmargin=1.4em,topsep=2pt,itemsep=1pt,parsep=0pt}

% tight float separation (CIKM: content must fit 10 pages incl. figures/tables)
\setlength{\textfloatsep}{8pt plus 4pt minus 3pt}
\setlength{\floatsep}{8pt plus 4pt minus 3pt}
\setlength{\intextsep}{8pt plus 4pt minus 3pt}
\setlength{\abovecaptionskip}{4pt}
\setlength{\belowcaptionskip}{2pt}

\newcommand{\sys}{\textsc{Qdcvr}}
\newcommand{\kbtool}[1]{\texttt{#1}}
\newcommand{\sig}[1]{$^{#1}$}

% ---------------------------------------------------------------------
% RED REVIEW MARKERS
% \need{...}      : inline gap the authors must fill before submission
% \needblock{...} : block-level gap (a table/figure/baseline that is missing)
% All red text is a to-do, never a claim. Remove every marker before
% submitting: a submitted PDF must contain no red text.
% ---------------------------------------------------------------------
\newcommand{\need}[1]{\textcolor{red}{\textbf{[\,#1\,]}}}
% \needblock renders as colour-shifted body text (not a box) so it participates
% in normal paragraph breaking — a minipage here produced 168pt overfull lines
% when placed directly after a \paragraph heading.
\newcommand{\needblock}[1]{%
  \par\medskip\noindent
  {\color{red}\small\textbf{[\,TO BE SUPPLIED\,]}\;#1}%
  \par\medskip}

% Allow TeX a little extra stretch so long technical compounds do not produce
% overfull lines in the narrow sigconf columns.
\setlength{\emergencystretch}{2.2em}
\tolerance=1200

% Generated macro facts (from the frozen snapshot; see make_assets.py).
\newcommand{\apiconsist}{480/480}
\newcommand{\apimismatch}{0}

\newcommand{\eFourRuns}{7}
\newcommand{\eFourFirsts}{6}
\newcommand{\eFourOursMean}{4.89}
\newcommand{\eFourLlmMean}{3.07}
\newcommand{\eFourRawMean}{3.39}

\newcommand{\qdcvrJhit}{8.48}
\newcommand{\qdcvrJmiss}{5.89}
\newcommand{\denseJhit}{9.18}
\newcommand{\denseJmiss}{8.25}
\newcommand{\qdcvrAbstain}{40}
\newcommand{\qdcvrJabs}{5.33}
\newcommand{\denseJabs}{8.09}


\begin{document}

\title{Content-Adjudicated, Domain-Scoped Retrieval for\\ Agent-Native Knowledge Bases}

\author{Anonymous Author(s)}
\affiliation{%
  \institution{Anonymous Institution}
  \country{}
}
\email{anonymous@example.org}

\begin{abstract}
Retrieval-augmented generation (RAG) pipelines over heterogeneous, multi-domain
corpora suffer two failure modes dense similarity alone cannot remove:
semantically adjacent but operationally irrelevant documents obtain high cosine
scores, and large bases monopolise the top-$k$ of a cross-base search. We
present \sys{}, a deployed knowledge-base platform whose retrieval is
\emph{content-adjudicated} and \emph{domain-scoped}: candidates are recalled by
a BM25\,$\rightarrow$\,vector cascade, then re-ranked by reading their actual
text against an interpretable $0$--$8$ rubric under an explicit
content-overrides-vector rule; an experience lifecycle distils operational
lessons into credibility-tiered entries exposed to agents through 94 Model
Context Protocol tools. Because a ranking metric is only a proxy for what
finally matters --- whether the returned text answers the question --- we
evaluate \emph{answer quality} with a three-agent protocol (one shared
answerer under a fixed evidence budget, an independent judge, a blind
middle-agent ranking) over an eight-system reproduced matrix on BEIR SciFact,
a nine-base HotpotQA partition, and an experience-synthesis comparison. The
boundary of the mechanism is precise: adjudication repairs rank-one errors
and yields the best top-rank answerability on the multi-domain benchmark
($0.580$); when its own retrieval ranks the gold document first, its answers
score \qdcvrJhit{} against the strongest baseline's \denseJhit{} ($0$--$10$
judge) --- the residual gap traces to its evidence-assembly policy, not to
adjudication itself. Distilled experiences beat raw retrieval and one-shot summarisation as
answering material in \eFourFirsts{} of \eFourRuns{} recorded runs. We also
report a negative oracle result for domain scoping, unstable experience
synthesis, and a post-freeze re-execution audit reproducing the eight-system
matrix byte-for-byte.
\end{abstract}
\begin{CCSXML}
<ccs2012>
<concept>
<concept_id>10002951.10003317.10003347</concept_id>
<concept_desc>Information systems~Retrieval models and ranking</concept_desc>
<concept_significance>500</concept_significance>
</concept>
<concept>
<concept_id>10002951.10003260.10003261</concept_id>
<concept_desc>Information systems~Question answering</concept_desc>
<concept_significance>300</concept_significance>
</concept>
<concept>
<concept_id>10010147.10010178.10010179</concept_id>
<concept_desc>Computing methodologies~Knowledge representation and reasoning</concept_desc>
<concept_significance>300</concept_significance>
</concept>
</ccs2012>
\end{CCSXML}

\ccsdesc[500]{Information systems~Retrieval models and ranking}
\ccsdesc[300]{Information systems~Question answering}
\ccsdesc[300]{Computing methodologies~Knowledge representation and reasoning}

\keywords{Retrieval-Augmented Generation; Content Verification; Knowledge
Management; Agent Systems; Model Context Protocol; Multi-Domain Retrieval}

\maketitle

% =====================================================================
\section{Introduction}\label{sec:intro}

Retrieval-augmented generation (RAG) has become the default way to ground large
language models in private knowledge \cite{lewis2020rag,gao2024ragsurvey}.
Almost all deployed pipelines share one architectural assumption: the document
that should be returned is the document whose embedding is closest to the query
embedding. Operating a production knowledge platform over a heterogeneous,
multi-domain corpus, we watched that assumption fail in two repeated,
measurable ways.

\noindent\textbf{Similarity is not usefulness.} Two documents can be close in
embedding space while only one of them answers the question. Query
\texttt{sf-002} of the BEIR SciFact benchmark asserts that \emph{``4-PBA
treatment decreases endoplasmic reticulum stress in response to general
endoplasmic reticulum stress markers''}; its gold document is PubMed id
32587939. Dense retrieval ranks a same-domain neighbour first and relegates
the gold document to rank two (Hit@1 $0$, MRR $0.5$) --- both documents live
in the same ``ER stress'' region of the space, but only one carries the
claim's evidence. The case is one of three such queries mined from the frozen
benchmark artefact (\texttt{sf-002}, \texttt{sf-017}, \texttt{sf-027}); two
of the three are repaired to Hit@1 $1.0$ (MRR $1.0$) when the adjudication
stage \emph{reads} the two candidates and inverts the order, while the third
keeps its miss --- the
concrete operation of ``content overrides vector''
(\S\ref{sec:adjudication}). A vector score is a joint statement about
the query and the document; it is not a statement that the document
\emph{contains} the answer.

\noindent\textbf{Big bases crowd out small ones.} In a multi-base deployment the
largest base dominates the candidate pool, and small but directly relevant bases
are pushed out. We measured this regime on a public corpus: 50 HotpotQA
questions whose official supporting documents are partitioned into nine topic
knowledge bases (\S\ref{sec:multidomain}); $38\%$ of queries have gold in two
of them.

Existing mitigations attack these problems from the wrong end: cross-encoder
re-rankers are opaque and add a model to the serving path; FLARE
\cite{jiang2023flare}, Self-RAG \cite{asai2023selfrag} and CRAG
\cite{yan2024crag} entangle the relevance judgement with generation or reduce
it to a scalar, so the operator still cannot see \emph{which} aspect of a
document failed.

\sys{} is a deployed knowledge-base platform built on a different principle:
\emph{recall cheaply, then adjudicate on content}. Candidates are recalled by a
two-stage cascade, then \emph{read} and scored against an explicit
three-dimension rubric; a content score overrides a high vector score. Queries
are \emph{domain-scoped} (routed to a bounded set of bases before comparison),
and operational knowledge is a first-class \emph{experience lifecycle}. Each
ingredient has a precedent --- domain routing for query complexity
\cite{jeong2024adaptiverag} and entity graphs
\cite{edge2024graphrag,huang2025hirag}; content verification in learned
re-rankers \cite{nogueira2019rerank} and evaluator models \cite{yan2024crag} ---
but none makes the retrieval decision \emph{legible} and \emph{authoritative}
at once. The contribution is the auditable combination, an answer-quality
account of where it helps, and an honest account of where it stops helping.

\paragraph{Contributions.}
\begin{itemize}
  \item \textbf{C1 (domain-scoped architecture).} A knowledge-base routing
  layer that selects a bounded working set before retrieval --- described as an
  architectural component, not a validated accuracy contribution: at this
  corpus scale an oracle restriction to the gold base does not beat search-all
  (\S\ref{sec:routing}), and we publish that negative result.
  \item \textbf{C2 (content adjudication).} A separable pipeline stage with an
  interpretable $0$--$8$ rubric and a \emph{content-overrides-vector} rule,
  with the regime boundary measured: it does not improve ranking on clean
  single-domain SciFact, and on multi-domain HotpotQA-in-9-bases it is
  \emph{significantly worse} than its own recall stage on full ranking
  ($-0.061$ nDCG@10, $p{=}0.018$) yet highest on top-rank answerability
  (answer@1 $0.580$ vs $0.480$--$0.520$) --- a rank-one repair mechanism for
  vocabulary-sharing queries, not a ranking improvement.
  \item \textbf{C3 (answer-quality attribution).} A three-agent evaluation
  protocol --- one shared answerer under a fixed 4{,}000-character evidence
  budget, an independent judge, and a blind middle agent --- plus a
  \emph{conditional attribution} analysis that separates retrieval-mechanism
  failures from evidence-assembly failures. It locates \sys{}'s answer-quality
  gap on SciFact in the excerpt policy, not in adjudication: with the gold
  ranked first, \sys{} answers score \qdcvrJhit{} vs the strongest baseline's
  \denseJhit{} (\S\ref{sec:main}).
  \item \textbf{C4 (experience lifecycle).} An end-to-end lifecycle for
  operational knowledge --- extraction, credibility tiering, temporal decay,
  agent-facing retrieval --- in which no freshly generated entry can present
  itself as authoritative; we report synthesis quality over the complete
  seven-run record of the material-quality comparison
  (\S\ref{sec:exp-analysis}).
  \item \textbf{C5 (agent-native surface).} The platform exposed as 94 Model
  Context Protocol tools \cite{anthropic2024mcp}, verified end-to-end by a
  committed external client: 26 of 26 checks across eight groups pass
  against the live deployment (\S\ref{sec:e2e}).
\end{itemize}

% =====================================================================
\section{Related Work}\label{sec:related}

\paragraph{Dense and hybrid retrieval.}
Dense passage retrieval \cite{karpukhin2020dpr} and contrastive retrievers
\cite{izacard2022contriever} optimise the recalled set; BM25
\cite{robertson2009bm25} and hybrid fusion remain strong, and no single
retriever dominates across domains (BEIR \cite{thakur2021beir}; BRIGHT
\cite{su2025bright} extends the stress test to reasoning-intensive queries).
\sys{} adopts the hybrid cascade as its recall stage
and does not claim it as a contribution: in our measurements it is an
acceleration device (\S\ref{sec:recall}).

\paragraph{Reranking and verification.}
Cross-encoder re-rankers \cite{nogueira2019rerank} and retrieve-rerank-generate
pipelines \cite{glass2022re2g} improve ordering with a learned scorer;
Self-RAG \cite{asai2023selfrag}, FLARE \cite{jiang2023flare} and CRAG
\cite{yan2024crag} interleave retrieval critique with generation, judging
relevance with a verdict entangled with the generator. \sys{} instead makes
adjudication a \emph{separate, auditable stage} whose output is a small set of
named rubric dimensions and whose verdict can override the recall ordering by
construction.

\paragraph{Knowledge organisation and graph RAG.}
Annotating text against an entity catalogue is an early form of this idea
\cite{tagme2010ferragina}; GraphRAG \cite{edge2024graphrag} and HiRAG
\cite{huang2025hirag} build entity or hierarchical structures for corpus-level
questions; Adaptive-RAG \cite{jeong2024adaptiverag} routes queries by predicted
complexity. Our routing layer is complementary but operates over
\emph{knowledge-base} granularity --- deciding \emph{which collections to
search} is classical federated IR
\cite{callan1995distributed,callan2000distributed} re-instantiated over
embedding-backed bases, here coupled with per-base content adjudication and an
auditable record of the routing decision.

\paragraph{Enterprise and agentic knowledge bases.}
AgenticRAG \cite{suresh2026agenticrag} moves grounding burden onto an agentic
harness over enterprise knowledge bases;
AutoKB \cite{sahay2025autokb} automates construction of structured domain
knowledge bases; DeepRead \cite{li2026deepread} exploits document structure to
improve agentic search --- the same industrial setting. They differ
in \emph{where the judgement lives}: they improve the agent's behaviour over a
retrieval result, whereas we make the retrieval stage itself adjudicate content
and expose that judgement for audit; such a harness on top of an adjudicating
retriever is the natural combination, and we have not evaluated it.

\paragraph{Agent tool-use for knowledge management.}
ReAct \cite{yao2022react} and Toolformer \cite{schick2023toolformer}
established tool-augmented LLM behaviour; the Model Context Protocol
\cite{anthropic2024mcp} standardises the interface, and MCP-Pyserini
\cite{ge2026mcppyserini} exposes IR resources as tools. \sys{} takes the same
direction at platform scale: the knowledge base is a tool surface the agent
operates, including ingestion, curation and experience management.

\paragraph{Experience and operational knowledge.}
Knowledge-management systems store documents; they rarely store
\emph{lessons}. We are not aware of a deployed RAG
platform that models experience with explicit credibility tiers and temporal
decay and then exposes it to an agent; \S\ref{sec:experience} describes our
attempt and \S\ref{sec:exp-analysis} where it falls short.

% =====================================================================
\section{System Architecture}\label{sec:system}

\sys{} is a deployed, cross-platform knowledge platform. This section describes
how it is built and, more importantly, \emph{why} it is built that way: three
architectural commitments constrain everything in \S\ref{sec:method} and
\S\ref{sec:experience}. Figure~\ref{fig:arch} shows the topology,
Table~\ref{tab:modules} a map of the
functional modules a reader can hold onto while the rest of the paper
evaluates them.

% Hand-drawn: end-to-end system architecture. Three services, five storage
% layers, one agent surface. Natural width ~ textwidth: no downscaling.
\begin{figure*}[t]
\centering
\resizebox{0.85\textwidth}{!}{%
\begin{tikzpicture}[
  font=\small,
  svc/.style={draw, rounded corners=2pt, align=center, fill=black!6,
              minimum height=11mm, inner sep=4pt},
  layer/.style={draw, rounded corners=1.5pt, align=center, fill=white,
                minimum height=9mm, minimum width=30mm, inner sep=4pt,
                font=\footnotesize},
  agent/.style={draw, dashed, rounded corners=2pt, align=center, fill=black!3,
                minimum height=10mm, inner sep=4pt},
  ar/.style={-{Stealth[length=2mm]}, semithick},
  wr/.style={-{Stealth[length=2mm]}, very thick},
  rd/.style={-{Stealth[length=2mm]}, densely dotted, thick},
  lbl/.style={font=\footnotesize\itshape, inner sep=2.5pt},
]

% ── clients ──────────────────────────────────────────────────────
\node[agent, minimum width=44mm] (human) at (0,0)
  {\textbf{Human} (browser)};
\node[agent, minimum width=44mm] (llm)   at (0,-1.9)
  {\textbf{LLM agent} (MCP client)};

% ── service tier ─────────────────────────────────────────────────
\node[svc, minimum width=48mm] (web)  at (5.6,0)
  {\textbf{Nuxt web layer} \texttt{:6789}\\[1pt]
   \footnotesize KB/document lifecycle, auth proxy};
\node[svc, minimum width=48mm] (mcp)  at (5.6,-1.9)
  {\textbf{kb-mcp server}\\[1pt]
   \footnotesize 94 MCP tools (stdio)};
\node[svc, minimum width=50mm] (api)  at (12.0,-0.95)
  {\textbf{FastAPI backend} \texttt{:8771}\\[1pt]
   \footnotesize retrieval, graph, experience, persona};
\node[svc, minimum width=44mm] (mineru) at (12.0,0.95)
  {\textbf{MinerU} OCR/parse\\[1pt]
   \footnotesize local or remote engine};

% ── storage layers ───────────────────────────────────────────────
\node[layer] (l1) at (0,-3.7)  {\textbf{L1} files on disk};
\node[layer] (l2) at (3.5,-3.7) {\textbf{L2} tree index\\ \scriptsize\ttfamily .tree-fs.json};
\node[layer] (l3) at (7.0,-3.7) {\textbf{L3} YAML catalogue\\ \scriptsize\ttfamily .kb.yml};
\node[layer] (l4) at (10.5,-3.7) {\textbf{L4} vectors\\ \scriptsize ChromaDB};
\node[layer] (l5) at (14.0,-3.7) {\textbf{L5} property graph\\ \scriptsize Neo4j};

% ── edges ────────────────────────────────────────────────────────
\draw[ar] (human) -- (web);
\draw[ar] (llm)   -- (mcp);
\draw[ar] (web) edge node[lbl]{HTTP} (api);
\draw[ar] (mcp) edge node[lbl, below]{HTTP} (api);
\draw[ar] (api) -- (mineru);
\draw[ar, dashed] (web.south) to[out=-90,in=130] (l3.north west);

% write path: API writes all five layers atomically (curves to L1/L2 dip
% through the corridor below the kb-mcp box so no line crosses node text)
\draw[wr] (api.south) -- node[lbl, right, pos=0.45]{single write path} (l4.north);
\draw[wr] (api.south west) to[out=-120,in=55] (l2.north east);
\draw[wr] (api.south west) to[out=-135,in=40] (l1.north east);
\draw[wr] (api.south east) to[out=-60,in=80] (l5.north);

% read path: retrieval reads indices directly, bypassing the API
\draw[rd] (mcp.south east) to[out=-30,in=170] (l4.west);
\draw[rd] (mcp.south) to[out=-80,in=95] (l3.north west);

% ── annotation ───────────────────────────────────────────────────
\node[lbl, anchor=west, align=left] at (-1.6,-5.1)
  {Solid = control flow;\; \textbf{very thick} = the single atomic write path (all five layers,
   one logical transaction);\; dotted = the read path used by retrieval, which bypasses the service tier.};
\end{tikzpicture}}
\caption{System architecture. Three services partition responsibility: the web
layer owns knowledge-base and document lifecycle, the FastAPI backend owns
parsing, retrieval, graph and experience computation, and MinerU does OCR
(self-hosted or as a remote engine). Every mutation travels one write path
that updates all five storage layers inside a single logical transaction
(\S\ref{sec:sys-layers}); retrieval reads those layers directly, so search
adds no load to the service tier. The same capability set is published twice
--- as REST endpoints for the human UI and as 94 MCP tools for agents --- so a
procedure validated in one surface is executable in the other.}
\label{fig:arch}
\end{figure*}


\subsection{Three services, one capability set}\label{sec:sys-services}
The platform is split so that each service owns one kind of work.
A \textbf{web layer} (Nuxt, \texttt{:6789}) owns the knowledge-base and document
lifecycle --- bases, documents, tags, previews --- and is the browser's only
network peer: it authenticates callers against the backend's token service and
proxies every request, so no CORS surface or backend URL is ever exposed.
A \textbf{compute backend} (FastAPI, \texttt{:8771}) owns retrieval, vector
indexing, graph construction, experience management and persona training.
\textbf{MinerU} does OCR and layout analysis for PDFs and scans; since
v~0.9 it runs in a \emph{dual-mode} configuration --- a self-hosted engine on
an ephemeral self-picked port, with a hosted remote engine as fallback --- so a
deployment can never silently attach to a stale engine instance.

The split is not merely organisational: the capability set is published
\emph{twice} --- as REST endpoints and as 94 Model Context Protocol tools ---
so a procedure validated in one surface is executable in the other. This is
what makes the platform \emph{agent-native} rather than agent-accessible.

\subsection{Five storage layers, one write path}\label{sec:sys-layers}
A knowledge base is materialised simultaneously as (L1) files on disk, (L2) a tree
index, (L3) a per-base YAML catalogue, (L4) a vector collection, and (L5) a
property graph: each layer answers a different access shape. The redundancy is
also the platform's main consistency hazard, and the design answer is one rule:
\emph{every mutation travels one write path} that applies to all five layers
inside one logical transaction, after which the ingestion module verifies that
all five agree: membership accuracy and storage completeness are both $1.000$
across two corpora (\S\ref{sec:setup}, Metrics). The rule is
what makes the rest of the system safe to build on: an agent that creates a
document cannot leave the index and the catalogue disagreeing.

\subsection{Write--read asymmetry, hierarchy, and design principles}\label{sec:sys-asymmetry}
Writes go through the service tier so that consistency is enforced in one place.
Reads used by retrieval \emph{bypass the service tier entirely} and read the
on-disk indices directly --- the asymmetry is deliberate: it makes the cheap
stages of retrieval essentially free, which is what allows an expensive
adjudication stage to exist at all (\S\ref{sec:honest}). Knowledge bases nest
arbitrarily: a parent base is a container whose children carry the documents and
are addressable individually and in aggregate, and the hierarchy is the unit the
routing layer operates on (\S\ref{sec:routing}). Three design commitments recur
below. \emph{Legibility over accuracy}: when choosing between a better opaque
decision and a slightly worse explainable one, choose explainable --- hence
named rubric dimensions rather than a score. \emph{Earned trust}: no artefact
is authoritative on creation; credibility accrues from use and decays with
time. \emph{Honest absence}: when evidence is thin, the system says so rather
than presenting a weak result as strong.

\begin{table}[t]
\centering
\footnotesize
\setlength{\tabcolsep}{4pt}
\caption{Top: module map --- each module is exercised through its production
interface by the benchmark harness; the last column names where the paper
reports it. Bottom: platform scale, measured 2026-09-14 by
\texttt{82\_system\_scale.py}; corpus counts are per benchmark corpus.
[14 skills = knowledge-base-family skills (20 registered incl.\ other
families); one backend endpoint added since.]}
\label{tab:modules}
\begin{tabular}{@{}p{1.62cm}p{3.75cm}p{1.25cm}@{}}
\toprule
Module & Responsibility & Reported \\
\midrule
Ingestion & MinerU OCR/parse (local or remote engine), chunking, five-layer
atomic write & \S\ref{sec:setup} \\
Search & \sys{} seven-stage retrieval (\S\ref{sec:method}) &
\S\ref{sec:main} \\
Organise & duplicate detection, tag generation and cleanup, graph build and
search & \S\ref{sec:ops} \\
Experience & eleven-stage lifecycle, credibility tiers, decay, five-path
retrieval & \S\ref{sec:exp-analysis} \\
Persona & per-base persona definition, training data and voice management &
\S\ref{sec:e2e} \\
Agent surface & 94 MCP tools + 14 skills mirroring the REST capability set &
\S\ref{sec:agent} \\
\midrule
\multicolumn{3}{@{}l}{\emph{Scale (running deployment / this paper's corpora)}} \\
MCP tools / skills / engines & 94 / 14 / 14 & \\
Backend endpoints / web routes & 113 / 124 & \\
Demo corpus, 3 KBs (EN/ZH/JA) & 16 docs, 20 queries & \S\ref{sec:main} \\
BEIR SciFact (official qrels) & 148 docs, 30 queries & \S\ref{sec:main} \\
HotpotQA partition, 9 KBs & 454 docs, 50 queries & \S\ref{sec:multidomain} \\
SQuAD v1.1 subset & 8 docs, 16 queries & \S\ref{sec:honest} \\
\bottomrule
\end{tabular}
\end{table}

% =====================================================================
\section{\sys{}: Content-Adjudicated Retrieval}\label{sec:method}

\sys{} stands for \emph{Query-Driven Content-Verified Retrieval} --- the
``content verification'' of the name is the content-adjudication stage
(\S\ref{sec:adjudication}); the two terms are interchangeable below. The pipeline
has seven stages (Fig.~\ref{fig:pipeline} summarised in
Algorithm~\ref{alg:pipeline}); we describe the stages that carry the method's
claims and treat the rest briefly.

% Hand-drawn (not data-generated): the QDCVR pipeline and where the two
% mechanisms under evaluation sit inside it. Natural width ~ textwidth.
\begin{figure*}[t]
\centering
\resizebox{0.80\textwidth}{!}{%
\begin{tikzpicture}[
  font=\small,
  stage/.style={draw, rounded corners=2pt, minimum height=11mm, minimum width=21mm,
                align=center, fill=black!4, inner sep=3pt},
  core/.style={draw, rounded corners=2pt, minimum height=11mm, minimum width=21mm,
               align=center, fill=black!16, line width=1.0pt, inner sep=3pt},
  store/.style={draw, dashed, rounded corners=2pt, minimum height=8.5mm,
                align=center, fill=white, inner sep=4pt, font=\footnotesize},
  ar/.style={-{Stealth[length=2mm]}, thick},
]

% ── main pipeline ────────────────────────────────────────────────
\node[stage] (q)   at (0,0)     {Query};
\node[stage] (s0)  at (2.35,0)  {S0\\ intent +\\ rewrite};
\node[core]  (s1)  at (4.9,0)   {\textbf{S1}\\ domain\\ scoping};
\node[stage] (s2)  at (7.45,0)  {S2\\ two-stage\\ recall};
\node[stage] (s3)  at (10.0,0)  {S3\\ dedup\\ + $\tau$};
\node[core]  (s4)  at (12.55,0) {\textbf{S4}\\ content\\ adjudication};
\node[stage] (s5)  at (15.1,0)  {S5\\ tiering};
\node[stage] (s6)  at (17.65,0) {S6\\ answer +\\ blind spot};

\foreach \a/\b in {q/s0, s0/s1, s1/s2, s2/s3, s3/s4, s4/s5, s5/s6}
  \draw[ar] (\a) -- (\b);

% ── stores ───────────────────────────────────────────────────────
\node[store] (kb)  at (4.9,-1.9)  {knowledge bases (hierarchical)};
\node[store] (idx) at (10.0,-1.9) {BM25 + ChromaDB index (full-corpus chunks)};
\node[store] (exp) at (15.1,-1.9) {experience store ($P_0/P_1/P_2$)};

\draw[ar, dashed] (s1) -- (kb);
\draw[ar, dashed] (s2) -- (idx);
\draw[ar, dashed] (s4.south) to[out=-100,in=140] (idx.north east);
\draw[ar, dashed] (exp) -- (s5);

% ── read path annotation ─────────────────────────────────────────
\draw[ar, dotted, thick] (s4.south west) to[out=-115,in=35]
   node[above, font=\footnotesize, pos=0.35]
   {\texttt{kb\_doc\_read}} (idx.west);

% ── highlight the two mechanisms ─────────────────────────────────
\node[draw, rounded corners=3pt, fit=(s1), inner sep=2.2mm, line width=1.2pt] {};
\node[draw, rounded corners=3pt, fit=(s4), inner sep=2.2mm, line width=1.2pt] {};

% ── agent surface ────────────────────────────────────────────────
\node[store, minimum width=118mm] (mcp) at (8.8,-3.6)
  {every stage reachable through the agent-native surface: 94 MCP tools + 14 agent skills (adjudication reads \emph{text}, not embeddings)};
\draw[ar, dashed] (s6.south) to[out=-90,in=20] (mcp.north east);

\end{tikzpicture}}
\caption{The \sys{} pipeline. Two stages carry the method's claims:
\textbf{S1 domain scoping} bounds the search to a small set of knowledge bases
before any document is compared, and \textbf{S4 content adjudication} reads
candidate text and scores it $0$--$8$, with authority to override the recall
order (Eq.~\ref{eq:override}). Every stage is reachable through the same MCP
tool surface the agents use.}
\label{fig:pipeline}
\end{figure*}


\begin{algorithm}[t]
\footnotesize
\caption{\sys{} retrieval}\label{alg:pipeline}
\begin{algorithmic}[1]
\Require query $q$, candidate bases $\mathcal{K}$, budget $k$
\State $q' \leftarrow \textsc{Rewrite}(q)$ \Comment{Stage 0: intent + rewrite}
\State $\mathcal{K}^{*} \leftarrow \textsc{Route}(q', \mathcal{K}),\; |\mathcal{K}^{*}| \le 3$
       \Comment{Stage 1: domain scoping}
\State $C \leftarrow \textsc{TwoStageRecall}(q', \mathcal{K}^{*})$ \Comment{Stage 2}
\State $C \leftarrow \textsc{Dedup}(C);\; C \leftarrow \{c \in C : s(c) \ge \tau\}$
       \Comment{Stage 3}
\ForAll{$c \in C$}
  \State $r(c) \leftarrow \textsc{Adjudicate}(\textsc{Read}(c))$
         \Comment{Stage 4: 0--8 rubric}
\EndFor
\State $T \leftarrow \textsc{Tier}(C, r)$; \textbf{content overrides vector}
       \Comment{Stage 5}
\State \Return $\textsc{Synthesise}(T)$ with blind-spot report
       \Comment{Stage 6}
\end{algorithmic}
\end{algorithm}

\subsection{Stage 0: Intent and Query Rewriting}\label{sec:intent}
The incoming query is classified into one of five intents (factual, procedural,
troubleshooting, comparative, navigational) and rewritten from colloquial
surface form into a declarative statement plus keyword set. Intent changes the
retrieval bias: troubleshooting queries consult the experience store before the
document store, whereas factual queries go straight to the vector index. We
attach no number to this stage: an earlier draft quoted a rewriting gain from
the withdrawn benchmark family (\S\ref{sec:provenance}), and no re-measurement
exists on the committed harness.

\subsection{Stage 1: Domain Scoping}\label{sec:routing}
The routing layer selects a bounded working set of knowledge bases before any
document is retrieved. Each base carries a name, a description and a document
catalogue; the router scores the rewritten query against those descriptors and
keeps at most three bases:
\begin{equation}
\mathcal{K}^{*} = \operatorname*{top\text{-}3}_{K_i \in \mathcal{K}}\;
   \mathrm{Sim}\bigl(q',\, \mathrm{desc}(K_i)\bigr).
\end{equation}
The intent is that a query whose vocabulary is generic (``heat management'',
``fast charging'') is pinned to the domain the user actually meant, and the
semantic neighbourhoods of unrelated domains never enter the candidate pool.
The quantitative claim that scoping reduces false positives came from the
benchmark we withdrew in \S\ref{sec:provenance}, so we state the mechanism and
not a number; a partial check is available in the traceable data, where all 20
in-house queries were served from the scoped stage rather than falling back to
an unscoped search.

We close part of this gap with an oracle experiment on the multi-domain
benchmark (\S\ref{sec:multidomain}): restricting retrieval to the single gold
knowledge base (a perfect single-base router) or even taking the \emph{best}
per-gold-base result yields \emph{lower} quality than the deployed
always-search-all mode (two-stage nDCG@10 $0.772$ restricted vs $0.841$ global;
best-restricted $0.795$). At 454 documents, cross-base distraction does not
crowd out gold documents, and single-base restriction is structurally fatal for
the $38\%$ of queries whose gold spans two bases (Recall@2 capped at $0.5$
under any single-base router). So at this scale scoping has no gain to deliver
even under an oracle; its value, if any, must be sought where cross-base
distraction dominates.

\subsection{Stage 2: Two-Stage Recall}\label{sec:recall}
Recall is a BM25 stage followed by vector refinement and graph-neighbour
expansion: BM25 supplies a sparse candidate set, and the vector stage re-ranks
inside it, so the expensive dense comparison runs over tens rather than
thousands of chunks (on SciFact $1.58$\,s staged vs $0.12$\,s one-shot dense;
at most 40 stage-1 candidates per query). We deliberately do \emph{not} claim
two-stage retrieval as a contribution: it is an engineering accelerator, and
our own standard-benchmark results show it can \emph{lose} to one-shot dense
retrieval on clean corpora (Table~\ref{tab:deepread}: Hit@3 $0.833$ vs
$0.900$); we report it as infrastructure.

\subsection{Stage 3: Deduplication and Thresholding}
Candidates are deduplicated by content fingerprint and filtered by a similarity
threshold $\tau$ (default $0.35$). The threshold is a precision/recall dial, and
it is the component most likely to hurt recall: on SciFact it can drop a gold
document whose similarity sits just below $\tau$. We analyse this failure
explicitly in \S\ref{sec:honest}.

\subsection{Stage 4: Content Adjudication}\label{sec:adjudication}
This is the stage that distinguishes \sys{} from a re-ranking pipeline.
Each surviving candidate is \emph{read} --- the top of its text is loaded through
the same document tool an agent would use --- and scored at temperature $0$ on
three named dimensions totalling $0$--$8$: \emph{topic relevance} ($0$--$3$;
$3$ = directly about the query subject), \emph{scenario match} ($0$--$3$;
$3$ = addresses the query's actual problem), and \emph{answer evidence}
($0$--$2$; $2$ = carries citable data, steps or conclusions). Scores come from
the document text, not from its embedding.

The decision rule makes the adjudication authoritative:
\begin{equation}
\mathrm{decision}(d) =
\begin{cases}
\text{accept }(P_0) & c(d) \ge 6\\
\text{supplement }(P_1) & c(d) = 5\\
\text{discard} & c(d) \le 4
\end{cases}
\label{eq:override}
\end{equation}
independently of the vector similarity $s(d)$,
where $c(d)$ is the content score and $s(d)$ the vector similarity. A document
with $s(d)=0.95$ and $c(d)=3$ is discarded. This is the operational content of
``vectors are fast, content is accurate''.

Three properties follow. The verdict is \emph{interpretable} (which dimension
failed), \emph{auditable} (the read text and scores are recorded, so a result
can be justified after the fact), and \emph{separable} (adjudication can be
disabled --- exactly what the ablation does). The harness exercises the
production procedure (recall, thresholding, adjudication via the same document
tool an agent uses; three verification reads per query, depth ablated); tag
expansion and answer synthesis are not exercised, so our claims concern the
retrieval core. The cost is latency: the adjudicator reads on average
$3{,}939$ characters per query in-house and $4{,}266$ on SciFact, which
\S\ref{sec:honest} treats as a first-class result.

\subsection{Stage 5--6: Tiering and Blind-Spot Declaration}\label{sec:blindspot}
Surviving documents are tiered $P_0/P_1/P_2$; chunks under $50$ characters are
demoted regardless of score, because such chunks receive inflated similarity
from embedding artefacts. When the confirmed set draws on fewer than two
knowledge bases the system declares coverage limited rather than presenting a
thin result as complete --- a design commitment, not a UI detail.

Two error directions are measurable. The rule is \emph{advisory by
construction}: it appends a coverage disclaimer and never removes a result, so
withholding-a-legitimate-result cannot occur in the measured pipeline (every
benchmark query returned results for every method). The
\emph{missed-declaration} direction is measurable on the multi-domain benchmark
(\S\ref{sec:multidomain}), where $38\%$ of queries need cross-base coverage:
taking ``the top-$5$ represents fewer gold bases than the gold set spans'' as
``should carry the disclaimer'', the missed rate is $26.3\%$ for two-stage,
dense and adjudicated, $63.2\%$ for BM25 --- the sparse path most needs the
declaration, as the rule's motivation predicts. The false-alarm direction is
decided by the answer-synthesis stage, which the harness bypasses; we state
that gap rather than infer a number.

% =====================================================================
\section{Experience Lifecycle}\label{sec:experience}

Operational knowledge --- what \emph{worked}, not merely what is known --- is
what agents need when a task fails, and most knowledge platforms do not model
it. \sys{} treats experience as a first-class artefact with its own lifecycle,
storage, credibility model and retrieval path.

\subsection{Lifecycle}\label{sec:lifecycle}
Experiences move through eleven stages: \emph{signal capture}, \emph{candidate
extraction}, \emph{draft synthesis}, \emph{review/approval},
\emph{publication}, \emph{vector indexing}, \emph{retrieval}, \emph{application
feedback}, \emph{review scoring}, \emph{credibility decay}, and \emph{health
monitoring}. Extraction runs in a deterministic heuristic mode or, via the
engine registry (14 agent engines), an LLM --- falling back silently to
heuristics when none is usable.

\subsection{Credibility Tiers}\label{sec:tiers}
A published experience is not equally trustworthy for every purpose, so
experience is tiered by combined evidence ($\mathrm{vec}$ = vector similarity,
$c$ = content score):
\begin{equation}
\mathrm{Tier}(e)=
\begin{cases}
P_0 & \mathrm{vec}\!\ge\!0.65 \;\wedge\; c\!\ge\!6 \;\wedge\; \mathrm{rating}\!\ge\!4 \;\wedge\; \mathrm{rev.}\!\ge\!1\\
P_1 & \mathrm{vec}\!\ge\!0.45 \;\wedge\; c\!\ge\!4\\
P_2 & \mathrm{vec}\!\ge\!0.35 \;\wedge\; c\!\ge\!3\\
\mathrm{discard} & \text{otherwise}
\end{cases}
\end{equation}
where $\mathrm{rating}$ is the entry's mean user rating on a $1$--$5$ scale
and $\mathrm{rev.}$ its count of independent reviews; the modifiers
\emph{downgrade} rather than upgrade: an entry disputed by
three or more reviews with mean rating below $2.0$ is capped at $P_2$; an entry
with no reviews and no applications is capped at $P_1$; automatically extracted
entries start unvetted at $P_1$. Trust must be earned through use; a freshly
generated experience can never present itself as authoritative.

\subsection{Temporal Decay}
Operational knowledge expires. Entries that are unverified and unapplied for
more than 30 days are demoted, and a scheduled decay pass re-evaluates the
store; decay is deliberately conservative --- it demotes, it never deletes ---
so a stale-but-historically-useful entry remains retrievable at a lower tier.
We evaluate the rule as a sensitivity analysis: the deployment is younger than
the window, so sweeping $7/14/30/90$ days the demoted set is empty at every
window -- correctly, since no entry is seven days old -- and no longitudinal
validation of the 30-day threshold is presented.

\subsection{Multi-Path Retrieval}
Experience is recalled along five parallel paths --- vector, keyword,
scenario, tag and quality-feedback --- fused and deduplicated: operational
queries often name a \emph{symptom} rather than a topic, sharing scenario and
tags instead of entry vocabulary. We evaluate the fusion on a purpose-built
set: eight operational experiences seeded through the platform's own write
path and twelve frozen queries --- eight \emph{symptom-style}, sharing no
vocabulary with the target entry's title, four \emph{vocabulary-style},
sharing solution words. Because the deployed endpoint exposes no per-path
switches, the fusion is re-implemented at harness level from the service's
published scoring rule; a top-3 agreement check against the live endpoint
matches $1$ of $6$ spot checks, so the numbers characterise the fusion
design, not the shipped endpoint's end-to-end behaviour.

\emph{Results.} The full five-path fusion is perfect on the twelve queries
(Hit@1 $=$ MRR $= 1.000$), and it is the only perfect configuration.
Leave-one-out: removing quality-feedback costs most (Hit@1 $0.833$), removing
keyword or scenario costs $0.083$ each, removing vector or tag none. Single
paths are all imperfect and fail on \emph{disjoint} queries (vector-only
$0.833$ fails the symptom-style queries with no lexical overlap; keyword-only
$0.917$ fails paraphrases; tag-only $0.417$; scenario-only $0.750$;
quality-only $0.167$) --- the complementarity the five-path design assumes.

% =====================================================================
\section{Agent-Native Interface}\label{sec:agent}

The platform is designed to be operated by an agent rather than browsed by a
human: it exposes \textbf{94 Model Context Protocol tools} covering
knowledge-base CRUD, document ingestion, parsing, tagging, search, vector
indexing, graph construction, experience management and persona training, plus
14 agent skills describing \emph{procedures} over those tools. Writes are
atomic and singular (a document creation does not silently index or tag it),
and the tool layer is the \emph{same} layer the human UI uses. Persona
training maintains a per-base persona definition with curated training data
and managed voice to steer generation style; the E2E client exercises this
end-to-end through its persona group (\S\ref{sec:e2e}).
End-to-end
validation of this surface: \S\ref{sec:e2e}.

% =====================================================================
\section{Evaluation}\label{sec:eval}

Retrieval metrics are proxies; the product of this platform is \emph{answered
questions}, so the evaluation combines standard benchmarks with an
\emph{answer-quality protocol} (\S\ref{sec:systems}) and answers six research
questions -- RQ1 does domain-scoped routing reduce cross-domain false
positives; RQ2 does content adjudication change ranking and answer quality,
and where does it pay; RQ3 do the individual components contribute; RQ4 does
the pipeline degrade on public, non-self-authored benchmarks; RQ5 what quality
does experience synthesis reach as \emph{answering material}; RQ6 is the
agent-facing surface correct end-to-end. RQ1 is answered by the oracle
experiment of \S\ref{sec:routing}: no gain even under an oracle at this scale.

\subsection{Benchmarks and Evaluation Protocol}\label{sec:setup}
\paragraph{Corpora.} Four corpora are used, each ingested through the production
pipeline and each reported with its own document count (Table~\ref{tab:modules}).
\emph{In-house demo corpus}: three knowledge bases (\texttt{KB-Demo-EN/ZH/JA})
holding 16 documents (15 Markdown in English, Chinese and Japanese plus one PDF
parsed through the production MinerU path), embedded with BGE-M3 (1024-d)
\cite{chen2024bgem3} and indexed in ChromaDB with a jieba BM25 keyword index;
the 20 bilingual questions of the in-house benchmark have their gold pages
inside these documents.
\emph{Public retrieval}: the BEIR SciFact collection \cite{thakur2021beir} ---
148 articles ingested through the production pipeline, evaluated on 30 official
queries against the official relevance judgements --- plus 16 SQuAD v1.1
questions. \emph{Multi-domain public}: 50 HotpotQA dev questions
\cite{yang2018hotpotqa} whose official supporting and distractor documents are
partitioned into nine topic knowledge bases by a deterministic keyword rule
(454 documents; $38\%$ of queries have gold in two bases); this benchmark
replaces a withdrawn internal 50-query set (\S\ref{sec:provenance}).

\paragraph{Metrics and functional coverage.} Ranking uses the standard
retrieval battery -- Hit@$k$, Recall@$k$, nDCG@10, Precision@5, MRR over
official qrels -- plus claim-word coverage. Beyond retrieval the benchmark
exercises every platform function through its production interface: parsing and
ingestion (Module~A: parse, ingest, membership and storage completeness all
$1.000$ across two corpora; self-retrieval Hit@1 $0.938$), the full experience
lifecycle (\S\ref{sec:exp-analysis}), the organise functions
(E17, \S\ref{sec:ops}) -- and the complete agent-facing HTTP surface (26/26 checks
across eight groups, \S\ref{sec:e2e}) over the 94-tool MCP layer.

\subsection{Compared Systems and the Answer-Quality Protocol}\label{sec:systems}

\paragraph{Reproduced baselines.} Following the experimental design of
\cite{li2026deepread}, \sys{} is compared against seven systems reimplemented
over the \emph{same} SciFact corpus and the \emph{same} frozen queries: dense
RAG over fixed 800/400-token windows; the same dense retriever with a reranker
(first-stage 30 candidates re-ranked, top 10 kept); RAPTOR's collapsed tree
(nodes $\le$800 tokens, five levels, cluster top-5, top-10 nodes retrieved);
two variants of iterative retrieval--generation synergy (four rounds, top-6
per round, evidence refreshed or accumulated); Search-o1-style agentic search
(two chunks per search, turn cap 8); and the DeepRead locate-then-read agent
itself. Every deviation from the original settings (embedding surrogate, LLM
listwise reranker in the RankGPT lineage \cite{sun2023rankgpt} in place of an
unavailable cross-encoder, turn cap chosen at the paper's own observed mean)
is registered in the reproduction notes.
\sys{} runs through its production MCP tool chain: two-stage recall, Stage-3
thresholding, and \texttt{kb\_doc\_read} content adjudication with the
content-overrides-vector rule --- the deployed retrieval path, not a
simulation.

\paragraph{Answer-quality protocol.} To attribute differences to retrieval
alone, every system's evidence is answered by \emph{one} shared agent under an
identical 4{,}000-character evidence budget and a frozen prompt; a second,
independent agent (fresh process, no shared context) grades each answer 0--10
after receiving the gold document; a third, \emph{middle agent} ranks all
eight anonymised answers per question. Prompts and channels are frozen, so
HTTP-driven and offline runs are byte-identical (\apiconsist{} checks,
Table~\ref{tab:deepread}, right columns).

\paragraph{Suite channels.} All suite channels read through the same tool
layer and index; only ranking differs: \emph{BM25} (sparse stage-1),
\emph{Two-stage} (+ vector refinement and thresholding, no adjudication),
\emph{Dense} (one-shot vector), and \emph{\sys{}} (+ content adjudication
and, where noted, domain scoping). Deterministic channels are free of
randomness (repeats bit-identical apart from latency); LLM channels run at
temperature $0$, reported as means with a $\pm1.5$ run-to-run tolerance on
judge means.

\paragraph{Implementation and hyper-parameters.} All configurations share one
index and one tool layer. Embeddings are BGE-M3 (1024-d, L2-normalised)
\cite{chen2024bgem3} from HuggingFace snapshot
\texttt{5617a9f6} via sentence-transformers $5.6.1$; chunks are $500$ characters
with $50$ overlap; the vector store is ChromaDB $1.5.9$ (persistent HNSW,
library-default $M{=}16$, efConstruction${=}100$, cosine); the keyword index is
jieba $0.42.1$ BM25 ($k_1{=}1.5$, $b{=}0.75$) over a $12{,}000$-character
window. Retrieval defaults are vector top-$k=10$, two-stage $k{=}40/10$,
threshold $\tau=0.35$, three verification reads; fusion is $0.5$ keyword and
$0.5$ graph-neighbour (depth $1$). The platform default for stage-1/2 is
$20/5$; the benchmark uses $40/10$, recorded in the frozen snapshot manifest.

\paragraph{Statistics.} The protocol is committed and executed
(\texttt{31\_stats.py}); every comparison names its numbers. Pairwise
comparisons use the \emph{Wilcoxon signed-rank
test} with a per-comparison $95\%$ bootstrap confidence interval ($B{=}2000$,
seed $0$) and \emph{Holm--Bonferroni} correction across the family. On the
multi-domain benchmark the corrected results are significant: adjudication
beats BM25 by $+0.241$ nDCG@10 ($p_{\mathrm{Holm}}{=}2\times10^{-5}$) and is
\emph{beaten} by its own two-stage recall stage ($-0.061$, $p{=}0.018$) and by
dense ($-0.088$, $p{=}0.018$). On SciFact ($n{=}30$) no component-level
difference survives correction (all $|\Delta| \le 0.023$): reported as
descriptive contributions with intervals. Deterministic retrieval repeats are
bit-identical, so we additionally ran a full reset-and-re-ingest reproduction
--- deployment wiped, all four corpora re-ingested, every benchmark re-run:
\textbf{every deterministic metric reproduced exactly}. The LLM channel
(experience synthesis) is inherently variable; we report nine state-reset runs
and analyse the variance as a finding (\S\ref{sec:exp-analysis}).

\subsection{Main Results}\label{sec:main}

\begin{table*}[t]
\centering
\footnotesize
\setlength{\tabcolsep}{3.6pt}
\caption{Main comparison on the BEIR SciFact corpus (148 documents ingested
through the production pipeline; 30 official queries, official relevance
judgements). All eight systems retrieve over the \emph{same} corpus and answer
the \emph{same} frozen questions; answers are produced by one shared agent
under a uniform 4{,}000-character evidence budget, then graded 0--10 by an
independent agent that receives the gold document (Judge). A middle agent then
ranks the eight anonymised answers per question given the gold evidence: Mean
rank = average position (1 = best of 8); Wins = first places. The right three
columns come from an HTTP-driven rerun that is byte-identical to the offline
run (480/480 consistency checks). Setting reproduced from
\citet{li2026deepread} on a claim-verification corpus.}
\label{tab:deepread}
\begin{tabular}{@{}lccccc ccc@{}}
\toprule
& H@1 & H@5 & R@5 & nDCG@10 & MRR & Judge & Rank & Wins \\
\midrule
Dense RAG (chunk 800/400, top-10) & 0.733 & 0.900 & 0.900 & 0.834 & 0.811 & \textbf{8.93} & 4.63 & 2 \\
Dense RAG + reranker (30$\to$10) & \textbf{0.900} & 0.933 & 0.933 & 0.921 & 0.917 & 8.67 & \textbf{3.77} & 5 \\
ITRG (refresh, 4 rounds $\times$ top-6) & 0.833 & 0.933 & 0.933 & 0.896 & 0.883 & 8.30 & 4.53 & 3 \\
ITRG (refine, 4 rounds $\times$ top-6) & 0.733 & 0.900 & 0.900 & 0.845 & 0.816 & 8.53 & 4.17 & 3 \\
RAPTOR (collapsed tree, top-10) & 0.733 & 0.867 & 0.856 & 0.817 & 0.800 & 8.37 & 4.20 & 4 \\
Search-o1 (agentic, 2 chunks/turn, cap 8) & 0.800 & 0.900 & 0.844 & 0.811 & 0.839 & 8.70 & 4.60 & 3 \\
DeepRead (locate-then-read, cap 8) & 0.567 & 0.733 & 0.694 & 0.635 & 0.639 & 8.13 & 4.47 & 7 \\
\sys{} (two-stage + content adjudication) & 0.700 & 0.833 & 0.833 & 0.784 & 0.767 & 7.70 & 5.63 & 3 \\
\bottomrule
\end{tabular}

\vspace{2pt}
\parbox{\textwidth}{\scriptsize H@$k$ = Hit@$k$; R@5 = Recall@5; nDCG@10 =
normalised discounted cumulative gain at 10 (official qrels). Judge grades are
LLM-as-judge with the gold evidence injected; a multi-grader robustness check
was not possible in our environment (single agent channel) and is stated as a
limitation. Deviations from the original settings are registered in the
reproduction notes (reranker surrogate, turn cap 8, 4 chars/token).}
\end{table*}


\paragraph{The eight-system matrix.} Table~\ref{tab:deepread} is the headline
comparison: seven reimplemented baselines and \sys{}, identical corpus,
identical questions, one shared answering agent, one independent grader. Three
readings. (i) \emph{Ranking}: the reranked dense baseline is strongest
(Hit@1 $0.900$, nDCG@10 $0.921$) --- on an abstract corpus a cross-encoder's
reordering is hard to beat --- while \sys{} sits mid-pack (Hit@1 $0.700$,
nDCG@10 $0.784$; the matrix runs \sys{} scoped to the corpus base, whereas the
suite-channel protocol with a global balanced pool reaches Hit@1 $0.767$):
adjudication re-orders the top of the list rather than improving overall
ranking (\S\ref{sec:ablation}).
(ii) \emph{Agentic behaviour is not free money}: Search-o1 reaches Hit@1
$0.800$ with two chunks per turn, but the DeepRead locate-then-read agent --
the strongest system in its own paper on long structured documents -- is
weakest here (Hit@1 $0.567$): abstracts carry no hierarchy for it to exploit,
so its extra turns buy context the corpus does not contain. We report this as
a corpus-boundary finding, not a refutation. (iii) \emph{Grades compress}:
answer-faithfulness grading narrows the spread ($7.70$--$8.93$;
per-question distribution in Fig.~\ref{fig:results}, left) because a correctly
supported ``insufficient'' verdict earns most of its points; the middle-agent
ranking (Rank/Wins columns) separates the systems more plainly and, notably,
ranks DeepRead first on $7$ of $30$ questions --- faithful abstention is
rewarded even when retrieval misses.

\begin{figure*}[t]
\centering
\includegraphics[width=0.455\textwidth]{generated/figures/fig-judge.pdf}\hfill
\includegraphics[width=0.455\textwidth]{generated/figures/fig-latency.pdf}
\caption{Left: per-question judge scores (eight systems $\times$ 30 SciFact
questions; mean bars with values; $\times$ = verdict not ok). Right:
retrieval latency (log) vs nDCG@10; marker area = mean judge score. Producer:
\texttt{make\_assets.py} from the frozen \texttt{deepread\_matrix.json}.}
\label{fig:results}
\end{figure*}

\paragraph{Where the answer-quality gap lives (RQ2).}
\sys{} has the \emph{lowest} judge mean ($7.70$) of the eight systems, and
Table~\ref{tab:ansqual} locates the deficit by conditioning each system's
judge scores on its \emph{own} rank-one success --- separating what is
retrieved from what is handed to the answerer. \emph{When a gold document is
ranked first}, \sys{} answers score \qdcvrJhit{} against dense's \denseJhit{}
--- the residual $0.7$ is excerpt truncation, not retrieval failure: on
\texttt{sf-028} the gold paper \emph{is} the retrieved top-1, but the GATA3
detail sits below the head-of-document excerpt horizon and the agent
cautiously declares insufficiency (judge $5$); on \texttt{sf-016}, whose
evidence fits the excerpt, the answer is graded $10$ with verdict, citation
and paraphrase matching the gold text. \emph{When rank one is missed},
\sys{} scores \qdcvrJmiss{} against dense's \denseJmiss{}: dense misses
``softly'' --- its ten excerpts ($67\%$ claim-word coverage) usually still
contain the gold --- whereas \sys{}'s verification reads keep $2$ excerpts
($60\%$ coverage). Abstention quality follows the same lever: \sys{} declares
insufficient evidence on \qdcvrAbstain\% of questions (dense $37\%$) but
scores \qdcvrJabs{} on them against dense's \denseJabs{}. The motivating
query \texttt{sf-002} shows the coupling in both directions: adjudication
\emph{repairs} the rank-one error (\S\ref{sec:intro}), and the answer then
declares insufficiency anyway (judge $7$ vs dense's $9$) --- the repaired
document won the ranking but its claim-bearing sentence never entered the
4{,}000-character budget. The gap is a property of the deployed
\emph{evidence-assembly policy}, not of adjudication itself; raising the
excerpt count is a configuration change, and \S\ref{sec:honest} quantifies
the same gap from the ranking side.

\begin{table}[t]
\centering\footnotesize
\setlength{\tabcolsep}{3.0pt}
\caption{Where the answer-quality gap lives (same 30 questions and judge
scores as Table~\ref{tab:deepread}): the judge mean conditioned on whether
the system's \emph{own} retrieval ranked a gold document first, the number
of rank-one misses, the share of answers that declare insufficient
evidence, and the judge mean on those abstentions.}
\label{tab:ansqual}
\begin{tabular}{@{}lccccc@{}}
\toprule
 & \multicolumn{2}{c}{Judge mean} & & Abstain & Judge mean \\
\cmidrule(lr){2-3}
System & gold@1 & no gold@1 & $n_{\text{miss}}$ & answers & on abstain \\
\midrule
Dense RAG & \textbf{9.18} & \textbf{8.25} & 8 & 37\% & \textbf{8.09} \\
Dense RAG + rerank & 8.74 & 8.00 & 3 & 37\% & 7.64 \\
ITRG (refresh) & 8.72 & 6.20 & 5 & 43\% & 7.23 \\
ITRG (refine) & 8.77 & 7.88 & 8 & 40\% & 7.58 \\
RAPTOR & 9.14 & 6.25 & 8 & 37\% & 7.09 \\
Search-o1 & 9.08 & 7.17 & 6 & 37\% & 8.00 \\
DeepRead & 9.06 & 6.92 & 13 & 47\% & 6.86 \\
\sys{} (ours) & 8.48 & 5.89 & 9 & 40\% & 5.33 \\
\bottomrule
\end{tabular}

\vspace{2pt}
\parbox{\columnwidth}{\scriptsize Producer: \texttt{make\_assets.py} from the
frozen \texttt{deepread\_matrix.json} (per-question judge, retrieval and
verdict rows).}
\end{table}


The suite channels (rows 1--3 of Table~\ref{tab:deepread}) and the in-house
demo corpus tell the same story from the ranking side. On the
\textbf{public benchmark} (30 official queries, official qrels, four methods
through one tool layer) the result is mixed and we report it as such: dense
retrieval has the best Hit@3 ($0.900$ vs $0.833$), Recall@5 and nDCG@10
($0.834$ vs $0.809$); BM25 has the best MRR ($0.839$); the content-verified
configuration wins only Hit@1 ($0.767$ vs $0.733$). On the \textbf{in-house
demo corpus} (20 queries over three knowledge bases, 16 documents) its only
lead is P@5 ($0.210$ vs $0.200$); it trails dense on Hit@1 ($0.800$ vs
$0.900$), Recall@5, and MRR ($0.863$ vs $0.931$), at $1.33$\,s against
$0.081$\,s. Verification \emph{re-orders the top of the list} rather than
improving ranking overall; whether it helps in proportion to cross-domain
distractors is tested next (\S\ref{sec:multidomain}).

\subsubsection{Withheld results and the provenance audit}\label{sec:provenance}
An earlier draft led with a 50-query internal benchmark (false-positive rate
falling from $0.635$ to $0.348$). \textbf{We have removed it}: the audit
found artefacts with \emph{no producing script}, an ablation on eight queries
rather than fifty, and a false-positive rate defined so that abstention scored
zero. The corrective is procedural: \textbf{every table must name its producing
script}. The withdrawn benchmark stays withdrawn; the multi-domain benchmark
replaced it (\S\ref{sec:multidomain}).

\subsection{Fine-Grained Behaviour, Ablation, and Reproducibility}\label{sec:behaviour}

\paragraph{Retrieval behaviour.} The cost structure separates the systems by
how they spend their budget (per-query means from the matrix artefact;
Fig.~\ref{fig:results}, right plots latency against ranking quality). Dense
runs $0.4$\,s and gathers ten excerpts covering $67\%$ of claim words; the
reranked variant spends $18.1$\,s, almost all in its single reranking call;
\sys{} spends $2.0$\,s and $3$ verification reads, keeping $2$ excerpts at
$60\%$ coverage; the iterative family spends $3$ agent calls per query to
reformulate ($60$--$64\%$); the agentic systems turn $2.7$--$3.7$ times at
$91$--$111$\,s ($44$--$53\%$). DeepRead gathers the least coverage of the
eight systems ($44\%$), and Search-o1, at $53\%$, still ranks mid-table --
reformulation buys relevance lexical overlap does not measure. The latency axis spans two orders of magnitude for a $\le 0.09$
nDCG@10 span; no method is simultaneously cheapest and best.

\subsubsection{Component ablation (RQ3)}\label{sec:ablation}
The component study is run on the \emph{same} thirty SciFact queries and the
same official qrels as the main table --- the query-set mismatch that made an
earlier component study unreportable has been removed. Every variant shares one
candidate pool per query: a single two-stage call whose stage-2 list supplies
the deployed configuration, and a deeper stage-2 call ($k{=}40$) that supplies
the threshold sweep. Table~\ref{tab:ablation} reports the committed harness
output (\texttt{30\_ablation.py}; within-session double-run is bit-identical).

\begin{table}[t]
\centering\small
\caption{Component ablation on BEIR SciFact (30 queries, official qrels), same
query set as Table~\ref{tab:deepread}. ``Deep'' variants re-rank a $k{=}40$ pool,
because the deployed top-10 is already pre-filtered by the Stage-3 threshold
(lowest returned score $0.379$), making a shallow-pool sweep a no-op --- itself
a finding; omitted rows in the artefact. Producer: \texttt{30\_ablation.py}.}
\label{tab:ablation}
\begin{tabular}{@{}lcccc@{}}
\toprule
Variant & Hit@1 & nDCG@10 & P@5 & MRR \\
\midrule
\sys{} full (deployed, $\tau{=}0.35$)  & \textbf{0.767} & 0.809 & 0.213 & 0.800 \\
\; verification $k{=}1$               & 0.733 & 0.796 & 0.213 & 0.783 \\
\; no verification (= Two-stage)      & 0.733 & 0.796 & 0.213 & 0.783 \\
\midrule
deep pool $k{=}40$, $\tau{=}0.20/0.35$/off & \textbf{0.767} & \textbf{0.832} & 0.213 & 0.810 \\
deep pool $k{=}40$, $\tau{=}0.50$     & \textbf{0.767} & 0.821 & 0.213 & 0.806 \\
\bottomrule
\end{tabular}
\end{table}

Three readings. (i) \emph{Content verification is load-bearing}: removing it
costs $0.012$ nDCG@10 and $0.017$ MRR, and shrinking it to one read reproduces
the no-verification numbers --- the second and third reads do the work.
(ii) \emph{Stage-3 thresholding is measurably harmful here, as
\S\ref{sec:honest} argues}: re-ranking the deeper $k{=}40$ pool at the same
threshold recovers $+0.023$ nDCG@10 over the deployed configuration (gold
documents the shallow pool never returned), while $\tau{=}0.50$ costs $0.011$
--- the deployed pool size, not the threshold value, loses gold documents.
(iii) Document dedup shows no incremental effect on this corpus, and no
component difference survives the multiplicity correction of
\S\ref{sec:eval}-Statistics ($|\Delta|\le0.023$): the table is a descriptive
decomposition, intervals in the artefact. The $2\times2$
scoping$\times$adjudication factorial remains open: the adjudication axis is
tested on the multi-domain benchmark (\S\ref{sec:multidomain}), the routing
axis is not (\S\ref{sec:routing}).
\subsubsection{RQ4: the honest reading on public benchmarks}\label{sec:honest}
SciFact is the evaluation least favourable to us and therefore the most
informative (Table~\ref{tab:deepread}): \textbf{content adjudication does not
improve ranking on a clean, single-domain benchmark}. The gold document
usually names the query entity, so vector similarity is already a good proxy
for usefulness, few cross-domain distractors exist for adjudication to remove,
and the Stage-3 threshold can drop a gold document whose similarity falls just
below $\tau$ -- part of the Hit@3 gap.

The distractor hypothesis --- that adjudication helps in proportion to
cross-domain distractors --- was the paper's central explanatory claim; the
figures that originally supported it came from the withdrawn family
(\S\ref{sec:provenance}). The multi-domain benchmark now tests it directly:
adjudication still repairs rank-1 errors on vocabulary-sharing queries (best
answer@1 of all methods; McNemar n.s., \S\ref{sec:multidomain}), but on
multi-hop questions with little lexical overlap it \emph{mis}-ranks and is
significantly worse than its own recall stage ($-0.061$ nDCG@10, $p{=}0.018$);
an overlap trigger does not recover this. Verification promotes documents
whose \emph{text} matches the question's terms; when none does, its
re-ordering is noise. Per-method coverage is reported as gold recall
(Table~\ref{tab:hotpot}) with the blind-spot analysis (\S\ref{sec:blindspot}),
replacing a false-positive rate that abstention could game.

\subsubsection{A multi-domain public benchmark: partitioning HotpotQA into
knowledge bases}\label{sec:multidomain}
Earlier multi-domain results rested on author-built data; we close that gap
with public data and official annotations: 50
questions sampled deterministically from the HotpotQA dev-distractor split
\cite{yang2018hotpotqa}, whose official supporting \emph{and distractor}
documents we ingest through the production pipeline and partition into nine
topic knowledge bases by a deterministic keyword-overlap rule (454 documents;
mapping frozen in the run manifest). Because a question's two gold documents
usually sit in different topics, $38\%$ of queries have gold in \emph{two}
bases --- cross-base coverage is a property of the official gold standard, not
of our partition. All four methods run through one tool layer with global
balanced pools; no routing is applied, so this benchmark isolates the
\emph{adjudication} axis of the factorial.

\begin{table}[t]
\centering\small
\caption{Multi-domain public benchmark: 50 HotpotQA dev questions, 454
supporting+distractor documents in 9 topic KBs, official supporting documents
as gold, answer-level check = gold answer string occurs in the read content of
the top-$k$ documents. Producer: \texttt{60\_hotpot\_build.py},
\texttt{61\_hotpot\_eval.py}.}
\label{tab:hotpot}
\begin{tabular}{@{}lcccccc@{}}
\toprule
Method & Hit@5 & R@2 & nDCG@10 & P@5 & MRR & ans@1 \\
\midrule
BM25       & 0.860 & 0.470 & 0.538 & 0.200 & 0.776 & 0.520 \\
Two-stage  & 0.960 & 0.720 & 0.841 & 0.344 & 0.930 & 0.500 \\
Dense      & \textbf{1.000} & \textbf{0.740} & \textbf{0.867} & \textbf{0.360} & \textbf{0.937} & 0.480 \\
\sys{}     & 0.960 & 0.620 & 0.779 & 0.344 & 0.820 & \textbf{0.580} \\
\bottomrule
\end{tabular}
\end{table}

Table~\ref{tab:hotpot} is the paper's most informative table: all methods are
unsaturated \emph{and} the corpus is multi-domain public data. With
Holm-corrected Wilcoxon tests over paired per-query nDCG@10: every staged
configuration beats BM25 by a wide margin (adjudicated $+0.241$,
$p_{\mathrm{Holm}}{=}2\times10^{-5}$); dense is strongest overall ($0.867$);
and --- the honest headline --- the adjudicated configuration is
\emph{significantly worse} than its own recall stage ($0.779$ vs $0.841$,
$p{=}0.018$) and than dense ($-0.088$, $p{=}0.018$): HotpotQA questions are
multi-hop and share little surface vocabulary with either candidate, so
word-matching verification often promotes the wrong document. The one place
adjudication leads is top-rank answerability (answer@1 $0.580$ vs
$0.480$--$0.520$; exact McNemar vs its recall stage $7/3$, $p{=}0.34$; vs
dense $p{=}0.23$; vs BM25 $p{=}0.58$ --- descriptive, not a tested effect).
With \S\ref{sec:honest} this yields the paper's bounded conclusion:
adjudication repairs rank-1 errors when query and evidence share vocabulary,
does not improve --- and can significantly hurt --- full ranking when they
do not.

\paragraph{In-house corpus and answerability.}
On the same demo corpus the split is language-dependent: on Chinese queries
\sys{} matches dense at Hit@1 $1.000$, whereas on English it is behind
($0.700$ vs $0.800$). Ranking metrics do not establish that the retrieved
text \emph{answers} the question, so we additionally read the full content of
the top-$k$ documents for every method and check it against the gold answer:
on SQuAD v1.1 ($n{=}16$) all four methods are saturated on ranking and
answer@1 is $0.938$ for the hybrid and vector channels ($0.875$ for BM25)
--- another honest null result: on a small, well-matched corpus the
adjudication stage has nothing to add.

\paragraph{Middle-agent ranking over the API surface.} The right three
columns of Table~\ref{tab:deepread} come from an HTTP-driven rerun: a
standalone client drives retrieval, answering, grading and a middle-agent
ranking of the eight anonymised answers on the public interface, and
reproduces the offline matrix byte-for-byte (\apiconsist{} consistency
checks) -- the benchmark is replayable through the same public surface an
external client would use. The middle agent's ordering (reranked dense best
at mean rank $3.77$) corroborates the qrels-based ranking, and answer quality
and retrieval quality separate exactly as the evidence-budget design intends.

\subsubsection{RQ5: Experience Synthesis as Answering Material}\label{sec:exp-analysis}
The lifecycle (\S\ref{sec:experience}) executes reliably: the scheduled
experience-synthesis pass (``meditation'' in platform terminology) succeeds,
tiers and decay behave as specified, multi-path retrieval
returns results. The \emph{quality} of automatically synthesised experience,
however, is unstable --- a negative result we report with its baselines and
its full run history. The measured outcome (artefact
\texttt{module\_c\_experience\_r2}, producer \texttt{03\_experience.py}):
10 experiences were produced for one base and none for the other (KB coverage
$0.50$), the mean judge score is $\mathbf{3.5/10}$, and only $37.5\%$ of
entries were judged grounded --- the two best were structurally sound but
carried ``no content body, parameters, or lessons''; the rest were ``empty
shells''. The architectural cause is identifiable: the synthesis stage
produces \emph{well-structured metadata with empty payloads}, which the
lifecycle caps at $P_1$ --- the failure degrades safely rather than injecting
noise into the authoritative tier. But a lifecycle that quarantines
low-quality output is not a system that produces good experience, and we do
not claim the latter.

\paragraph{Multi-run stability under state reset.}
An earlier freeze reported a different outcome; an audit hypothesised state
contamination. We made the reset part of the benchmark --- every round
deletes all published experiences and rejects all drafts, verifying zero
counts --- and have run nine rounds in three samples. Instability \emph{persists
under reset}, refuting contamination: the English base produced
$0,3,0$ / $3,0,0$ / $0,0,3$ entries, the Chinese base zero each time. Inside
the variance: when the agent \emph{does} produce, entries are judged
$7$--$9/10$ under the same strict judge that gave the heuristic table above
$3.5$, and the Chinese base's zero is the quality gate correctly refusing
encyclopedic content. The variance is the generative agent's produce/skip
judgement, not platform state; we report it rather than selecting a favourable
round.

\paragraph{Does distilled experience answer better? (RQ5).}
We compare three answering materials under an \emph{identical} judge prompt
(grounded 4 / structure 3 / reusable 3; four frozen operational queries ---
\texttt{x-01}, \texttt{x-02}, \texttt{x-04}, \texttt{x-08} --- sampled from
the eight-query operational set):
(i) \emph{no synthesis} --- the raw top-$3$ documents the two-stage tool
retrieves; (ii) \emph{one-shot LLM summary} --- the seven source documents
summarised into ``lessons'' by a single prompt; (iii) \emph{ours} --- the
pipeline-created experiences. Table~\ref{tab:e4} reports \emph{every recorded
run} of this comparison --- \eFourRuns{} in total, none omitted: the
distilled material ranks first in \eFourFirsts{} of \eFourRuns{} runs (mean
\eFourOursMean{} vs \eFourLlmMean{} one-shot and \eFourRawMean{} raw): the
seven-run \emph{mean} margins over both baselines ($1.82$ vs one-shot, $1.50$
vs raw) clear the $\sim$1.4 judge noise floor, while individual winning-run
margins range $0.25$--$2.75$ --- three of the six below the floor. The one
loss (R7) is the
instability above in miniature: that round's synthesis produced shell entries
(three of four sampled scores $0$), and the arm fell to last. When the
lifecycle distils real lessons, the result is the strongest of the three
materials; the open problem is making it distil reliably.

\begin{table}[t]
\centering\footnotesize
\setlength{\tabcolsep}{2.6pt}
\caption{Experience material as answering evidence: mean judge score (0--10,
identical judge prompt, eight frozen operational queries) over \textbf{all
seven recorded runs} of the dual-baseline comparison --- no round is omitted.
The pipeline arm ranks first in 6 of 7 runs;
in the exception (R7) its synthesis round produced shell entries, the
produce/skip instability of \S\ref{sec:exp-analysis}. Producers:
\texttt{40\_experience\_suite.py}, \texttt{42\_e4\_fix.py}.}
\label{tab:e4}
\begin{tabular}{@{}lccccccc|c@{}}
\toprule
Condition & R1 & R2 & R3 & R4 & R5 & R6 & R7 & Mean \\
\midrule
\sys{} (distilled) & \textbf{5.25} & \textbf{5.25} & \textbf{6.50} & \textbf{5.50} & \textbf{5.75} & \textbf{4.50} & 1.50 & 4.89 \\
One-shot summary & 3.00 & 2.25 & 3.50 & 3.00 & 3.75 & 2.50 & \textbf{3.50} & 3.07 \\
Raw documents & 4.00 & 2.50 & 3.75 & 4.50 & 2.00 & 4.25 & 2.75 & 3.39 \\
\bottomrule
\end{tabular}
\end{table}


\emph{Judge reliability:} across three independent runs of 24 calls scored
twice under two prompts sharing one rubric, agreement is moderate (exact
$22$--$25\%$, within one point $54$--$63\%$, $\kappa{=}0.22$--$0.40$ on the
binarised ``usable'' decision), so effects above $\sim$1.4 points survive;
this is an LLM--LLM agreement, not a human study. Producers:
\texttt{40\_experience\_suite.py},
\texttt{42\_e4\_fix.py}, \texttt{66\_judge\_agreement.py}; transcripts
archived.

\subsubsection{RQ6: End-to-End Verification of the Agent Surface}\label{sec:e2e}
We verify the externally reachable surface with a committed, independently
written client that never imports application code
(\texttt{80\_e2e\_surface.py}): authentication; knowledge-base and document
lifecycle; content-based retrieval including a staged read; graph service; the
experience lifecycle; persona; the engine registry; cleanup --- all over HTTP
as an agent would call them. The current run passes \textbf{26 of 26 checks
across eight groups} against the live deployment and writes a JSON artefact
carrying its own git and configuration fingerprint; it also surfaced two
contract defects, both fixed. This 26-check committed client is the paper's
frozen verification gate; the platform's external-API suite has since grown to
73 checks, a later expansion that does not replace it.

\subsubsection{Organise functions (E17)}\label{sec:ops}
A planted-truth probe (\texttt{26\_platform\_ops\_eval.py}) exercises the
\emph{organise} functions on a purpose-built nine-document base seeded with
two duplicate groups. Findings, reported rather than smoothed over: the
creation chain silently drops exact-duplicate documents (nine planted, seven
persist, no error signalled), so the duplicate detector sees only the
near-duplicate pair and identifies one of two groups; tag generation is noisy
on small documents ($2.6\%$ of 190 tags grounded in document content);
cleanup dry-run never touches a used tag; graph build succeeds; the catalog
mirrors exactly the seven persisted documents with no spurious entries (its
completeness flag is false precisely because the dropped duplicates leave
seven of the nine expected documents).

\subsubsection{Reproducibility of the reported numbers}
Every number is generated from frozen artefacts, never transcribed: the bytes
are archived with a SHA-256 manifest that the table generator verifies on every
run; each run writes to an immutable timestamped directory embedding its git
commit, configuration hash, seed and run id; every table names its producing
script (enforced by \texttt{provenance\_audit.py}).
After this manuscript was frozen we re-executed both benchmark
tracks from the committed pipelines and machine-compared every artefact
against the frozen bytes (grades within the $\pm1.5$ LLM tolerance; latency,
timestamps and run provenance excluded): the eight-system matrix reproduces
\textbf{9{,}837 of 9{,}837 fields byte-identically}, the platform-functions
track is aggregate-identical on all five artefact pairs, and two fresh
real-document runs agree on \textbf{47 of 48} gold-document positions
(97.9\%). The audit also fixed four pipeline defects and reports one boundary
finding as-is: the multi-domain benchmark's \emph{absolute} numbers depend on
the vector store's build-time state, while the \emph{ordering} reproduced in
every healthy run --- so \S\ref{sec:multidomain} holds on either set of
numbers. Full log: artefact.

% =====================================================================
\section{Discussion and Limitations}\label{sec:conclusion}

\paragraph{What we would claim, and what we would not.}
The defensible claims are narrower than when we began. Ingestion is lossless.
Content adjudication
\emph{re-orders} the top of a list rather than improving ranking overall ---
but the re-ordering is exactly what answerability needs: it is the best
rank-one answer producer on the multi-domain benchmark, and when its evidence
channel covers the gold document its answers are graded within a point of the
strongest baseline. The answer-quality gap that remains is located, measured,
and attributable to the excerpt policy --- a configuration, not the mechanism.
The motivating claim --- that scoping plus adjudication substantially reduces
false positives in heterogeneous corpora --- is \emph{refined}: the
multi-domain benchmark bounds the adjudication half, the scoping half remains
unmeasured at this scale.

\paragraph{Latency and evaluation scale.}
Content adjudication costs about $16\times$ in latency on the in-house corpus
($1.33$\,s vs $0.081$\,s); reads bypass the service tier and apply only to a
small candidate set --- not recommended in latency-critical single-domain
settings. The public evaluations are moderate (30 SciFact, 50 HotpotQA, 16
SQuAD queries; full-BEIR is the obvious next step).

\paragraph{Threats to validity.}
\emph{Construct}: the $0$--$8$ rubric and the experience judge are LLM
judgements; we mitigate with temperature $0$, frozen prompts and a dual-judge
agreement study (\S\ref{sec:exp-analysis}: $\kappa{=}0.22$--$0.40$) -- a proxy
for human agreement.
\emph{Internal}: the component study shares the thirty SciFact queries of
Table~\ref{tab:deepread}. \emph{External}: the in-house corpus is self-authored;
the HotpotQA partition lifts the single-domain restriction but is
author-chosen. \emph{Author-evaluation bias}: all benchmarks were designed, run
and scored by the system's authors; we mitigate with public gold standards and
unfavourable results. \emph{Coverage}: the benchmarks measure the retrieval
core; the rest of the platform is exercised for correctness, not quality.

\paragraph{Future work.}
The evidence-assembly policy is the highest-leverage repair: \sys{}'s
verification reads keep two excerpts where dense keeps ten, and
Table~\ref{tab:ansqual} shows the answer-quality cost of that choice. The
experience-synthesis instability is generative, not architectural: the repair
path is to stabilise the produce/skip decision. The Stage-3 threshold should be learned
per corpus and the stage-2 pool deepened: the pool size, not the threshold
value, loses gold documents. Adjudication should condition on query--evidence
lexical overlap; scoping must be tested at corpus sizes where cross-base
distraction dominates; and a human inter-annotator study for the judge rubric
remains the missing validation.

\paragraph{Conclusion.}
We described \sys{}, a deployed agent-native knowledge platform, and evaluated
its two central retrieval mechanisms and its answer quality on public
benchmarks --- BEIR SciFact with a reproduced eight-system matrix, a new
multi-domain HotpotQA partition, a bilingual in-house corpus, and a
state-resetting multi-round experience protocol. The picture is bounded:
content adjudication is a rank-one repair instrument for queries that share
vocabulary with their evidence, and its answers are competitive \emph{whenever
its evidence channel lets them be} --- the measured gap lives in excerpt
assembly, not in the mechanism; distilled experience, when the lifecycle
produces it, beats raw retrieval and one-shot summarisation as answering
material; ingestion is lossless; the lifecycle contains unstable generation
safely; and scoping shows no gain even under an oracle at this scale.
Negative results of this kind tell the next system where the mechanism stops
working.

% =====================================================================
\section*{GenAI Usage Disclosure}
Generative AI tools were used per the ACM Authorship Policy. \emph{Writing}:
language editing only. \emph{Code}: parts of the system and benchmark harness,
shipped with the artefact under the authors' responsibility.
\emph{Data}: no result was produced or modified by a generative model; every
number is read from a recorded artefact by deterministic scripts. One LLM is
used as an instrument: the experience judge (\S\ref{sec:exp-analysis}).
Research questions, study design and interpretation are the authors'.

\balance
\bibliographystyle{ACM-Reference-Format}
\bibliography{refs}

\end{document}
